\chapter{Судалгааны сэдвийн онол, өнөөгийн түвшин}
\label{ch:literature}

%======================================================================
% ЭНЭ ФАЙЛД БАЙГАА БҮХ ЗҮЙЛ БАГШИД ҮЗҮҮЛЭХ НООРОГТ ХЭВЛЭГДЭНЭ (бүтэн ноорогт ч мөн).
% Зөвхөн бүтэн ноорогт орох хэсэг тусдаа файлд: Chapters/full/ch2_*.tex (\fullinput мөрүүд).
%
% 2026.11.16 (2), хэрэглэгчийн шүүмжээр дахин бичсэн: бүлэг яриа танилтад хэт төвлөрсөн байсныг
% тэнцвэржүүлсэн — эхлээд системийн гол хэсэг (хайлтаар баяжуулсан хариулт, хайлтын аргууд),
% дараа нь яриа таних (богиносгосон; үйлчилгээний квот, нөхцөлийн догол мөр зөвхөн бүтэн ноорогт).
% 2026.11.16 (3), хэрэглэгч: мэдлэгийн мод бол RAG-ийн суурь — «Мэдлэгийн мод дээрх хайлт» хэсгийг
% бүтэн нооргоос багшид үзүүлэх ноорог руу шилжүүлсэн (эшлэл, тоо хуучин хэвээр).
% Шинэ эшлэл, шинэ тоо нэмээгүй. Өмнөх хувилбар: thesis/archive-20261116/.
% Whisper, Mengke, Lewis-ийг хэрэглэгч өөрөө уншсан (✅). Бусад эшлэлийг агент бүтэн уншиж
% papers/-д бүртгэсэн (papers/16, 24, 27); тоо, мета өгөгдлийг 2026.11.15-нд эх хуудаснаас тулгасан.
%======================================================================

Энэ бүлэгт асуултад баримтын сангаас хайж хариулах арга, мэдлэгийн модыг хайлтад ашиглах арга болон Монгол хэлний яриа танихтай холбоотой судалгааг авч үзэв.

\section{Хайлтаар баяжуулсан хариулт (RAG)}
\label{sec:lit-rag}

Хэлний том загвар сургалтын үедээ олон баримтыг жиндээ хадгалдаг. Lewis нар үүнийг \textit{жиндээ хадгалсан мэдлэг} (parametric memory) гэж нэрлэсэн \cite{lewis2020rag}. Ийм мэдлэгийг шинэчлэхийн тулд загварыг дахин сургах хэрэгтэй. Мөн загвар хариултаа аль эх сурвалжаас авснаа харуулж чаддаггүй, заримдаа бодит бус мэдээлэл зохиодог \cite{lewis2020rag}.

Хайлтаар баяжуулсан хариулт (RAG) нь загварт \textit{гадаад санах ой}, өөрөөр хэлбэл тусдаа хадгалсан бичвэрийн сан нэмдэг. Ийм систем хоёр хэсэгтэй \cite{lewis2020rag}:
\begin{enumerate}
    \item \textbf{Хайгч} (retriever) асуултад хамгийн ойр $K$ баримтыг сангаас олно. Асуулт ба баримтыг тоон вектор болгоод, векторын ойролцоо байдлаар эрэмбэлнэ.
    \item \textbf{Үүсгэгч} (generator) асуулт ба олдсон баримтыг хараад хариулт бичнэ.
\end{enumerate}

Нээлттэй асуулт-хариултын Natural Questions өгөгдөл дээр RAG-Sequence загвар яг таарах (exact match) хэмжүүрээр 44.5\% авсан бол зөвхөн жиндээ мэдлэгтэй T5-11B+SSM загвар 36.6\% авсан \cite{lewis2020rag}. Мөн баримтын санг 2016 оны Wikipedia-гаас 2018 оных болгож солиход 2018 оны улсын тэргүүний тухай асуултын зөв хариулт 4\%-иас 68\% болсон. Өөрөөр хэлбэл загварыг дахин сургахгүйгээр зөвхөн санг сольж мэдлэгийг шинэчилж болно \cite{lewis2020rag}. Сургуулийн мэдээлэл жил бүр өөрчлөгддөг тул энэ шинж чанар энэ ажилд чухал.

Хайгч хэсэгт хэд хэдэн аргыг хэрэглэдэг. BM25 нь асуултын үг баримтад хэдэн удаа орсныг тоолж оноо өгдөг бөгөөд ховор үгэнд их жин өгч, баримтын уртыг тооцдог \cite{robertson2009bm25}. Multilingual E5 нь олон хэлний бичвэрийг тоон вектор болгодог загвар бөгөөд үг нь таараагүй ч утга нь ойр бичвэрийг олоход ашиглагдана \cite{wang2024me5}. \acs{RRF} (Reciprocal Rank Fusion) нь хэд хэдэн хайлтын жагсаалтыг баримт бүрийн эзэлсэн байраар нь нэгтгэдэг арга юм \cite{cormack2009rrf}. Энэ ажлын хайлт эдгээр гурван аргыг хослуулсан (\ref{sec:rag-method}-р хэсэг).

Энэ ажилд RAG-ийн хайгч хэсгийг ашиглаж, хайлтаар олдсон эх сурвалжийн өгүүлбэрийг үгчлэн харуулна (\ref{sec:requirements}-р хэсэг). Lewis нарын дээрх үр дүн хариултыг хэлний загвараар бичүүлдэг системийнх тул энэ системийн хариултын чанарыг шууд батлахгүй.

\fullinput{Chapters/full/ch2_unanswerable}

\section{Мэдлэгийн мод дээрх хайлт}
\label{sec:lit-tree}

Мэдлэгийн мод гэж мэдээллийг сэдвээр нь шатлан ангилсан бүтцийг хэлнэ. Жишээ нь «Элсэлт $\rightarrow$ Төлбөр», «Байршил $\rightarrow$ Өрөө». Ийм модыг хайлтад хэд хэдэн аргаар ашиглаж болно.

\fullinput{Chapters/full/ch2_tree_literature}

Гараар байгуулсан модыг хайлтад ашиглах гурван хувилбарыг \tref{tab:tree_options}-д харьцуулав.

\begin{table}[H]
\centering
\caption{Мэдлэгийн модыг хайлтад ашиглах гурван хувилбар}
\label{tab:tree_options}
\begin{tabular}{|p{2.6cm}|p{5.7cm}|p{5.3cm}|}
\hline
\textbf{Хувилбар} & \textbf{Юу хийх} & \textbf{Эрсдэл} \\
\hline
Энгийн хайлт & Бүх эх сурвалжаас хайна. Мод зөвхөн хариулт олдохгүй үед хэнд, хаана хандахыг заахад ашиглагдана & Өөр сэдвийн төстэй хуудас дээгүүр гарч болно \\
\hline
Хүрээний нэмэлт оноо & Энгийн хайлт дээр нэмж асуулт аль хүрээнд хамаарахыг тодорхойлж, тэр хүрээний эх сурвалжид нэмэлт оноо өгнө. Бусдыг хасахгүй & Нэмэлт оноо хэт их бол буруу хүрээ давамгайлна \\
\hline
Салаагаар шүүх & Эхлээд тодорхойлсон салаандаа л хайна & Салааг буруу тодорхойлбол зөв эх сурвалж огт хайгдахгүй \\
\hline
\end{tabular}
\end{table}

Салаагаар шүүх хувилбарт салааг буруу тодорхойлбол зөв эх сурвалж огт хайгдахгүй. Иймд энэ ажилд хүрээний нэмэлт онооны хувилбарыг сонгосон: асуултыг 11 том хүрээний түвшинд ангилж, тэр хүрээний эх сурвалжид нэмэлт оноо өгнө, бусад эх сурвалжийг хасахгүй. Энэ нь модоор алхаж хайдаг арга биш, асуултыг ангилаад хайлтын оноонд зөөлөн нөлөөлдөг арга юм. Жишээ нь «элсэлтийн төлбөр хэд вэ» гэсэн асуултыг «Элсэлт» хүрээнд хамааруулж, элсэлтийн хуудсуудыг жагсаалтын дээд талд гаргана. Аргын хэрэгжилтийг \ref{sec:rag-method}-р хэсэгт тодорхойлов.

\section{Монгол хэлний яриа таних}
\label{sec:lit-asr}

Яриа таних (\acs{ASR}) гэдэг нь хүний ярианы аудиог текст болгох ажил юм. Хүн асуултаа дуугаар асуухад систем эхлээд яриаг текст болгож, дараа нь тэр текстээр хайна. Иймд яриа танилтын алдаа хайлтад шууд нөлөөлнө. Англи хэл дээрх нэг туршилтад асуултыг дуу үүсгэгчээр уншуулж, яриа танигчаар буцаан текст болгоод хайхад BM25 хайлтын MRR@10 үзүүлэлт (зөв баримт эхний 10 үр дүнгийн хэддүгээр байрт гарсны урвуу утгын дундаж) 18.67-оос 6.97 болж буурсан \cite{sidiropoulos2022asr}.

Яриа танилтын чанарыг ихэвчлэн үгийн алдааны хувиар (\acs{WER}) хэмждэг:
\begin{equation}
\mathrm{WER} = \frac{S + D + I}{N} \times 100\%
\label{eq:wer}
\end{equation}
Энд $S$ нь буруу таньсан үг, $D$ нь орхисон үг, $I$ нь илүү нэмсэн үгийн тоо, $N$ нь жишиг бичвэрийн (хүн сонсож бичсэн хувилбар) нийт үгийн тоо. Илүү нэмсэн үг олон бол WER 100\%-иас давж болно.

Whisper нь OpenAI-ийн нээлттэй яриа таних загвар юм \cite{radford2023whisper}. Үүнийг интернэтээс авсан 680 мянган цагийн аудио ба бичвэрээр сургасан бөгөөд үүний 117 мянган цаг нь англиас бусад 96 хэлнийх. Whisper large-v2 загварын Монгол хэлний WER FLEURS өгөгдөл дээр 110.5\%, Common Voice 9 дээр 128.4\% гарсан \cite{radford2023whisper}. % Хавсралтын Table 13, Table 11-тэй агент тулгасан (11/06).

Mengke нар Халх Монгол хэлний яриа танилтыг сайжруулсан \cite{mengke2025khalkha}. Common Voice 17-д Халх Монгол хэлний сургалтын өгөгдөл ердөө 3 цаг байдаг. Зохиогчид Өвөр Монголын Чахар аялгууны 147 цагийн өгөгдлийн бичвэрийг кирилл бичигт хөрвүүлж ашигласан. Whisper large-v3 загварыг шууд хэрэглэхэд WER 92.60\% байсан бол Халхын 3 цагийн өгөгдлөөр сургахад 43.02\%, эхлээд Чахарын, дараа нь Халхын өгөгдлөөр сургахад 31.45\% болсон \cite{mengke2025khalkha}. Энэ нь Whisper-ийг Монгол хэлэнд сайжруулж болохыг харуулсан. Гэхдээ үүнд өгөгдөл, хүчтэй график процессор (\acs{GPU}) хэрэгтэй.

Системд бэлэн ашиглаж болох хоёр боломж бий. Нэг нь Монгол хэлээр нэмэлт сургасан нээлттэй загвар юм. Тухайлбал Whisper small-д суурилсан \texttt{Cafet/\allowbreak whisper-small-\allowbreak mongolian} загвар интернэтгүйгээр, өөрийн компьютер дээр ажиллана \cite{cafet_whisper_small_mn}. Нөгөө нь онлайн үйлчилгээ юм. Google-ийн Gemini 3.5 Transcribe үйлчилгээ Монгол хэлийг (\texttt{mn-MN}) дэмждэг \cite{google_transcribe_model}. Аудио Google-ийн сервер рүү очдог тул интернэт хэрэгтэй. Үнэгүй хэрэглээтэй, төлбөртэй хувилбар нь минутад 0.005 ам.доллар орчим \cite{google_pricing}. Энэ үйлчилгээний Монгол хэлний алдааг хэмжсэн судалгаа энэ ажлын явцад хийсэн хайлтаар олдоогүй. Хоёр боломжийн гол ялгааг \tref{tab:asr_compare}-д харуулав.

\fullinput{Chapters/full/ch2_asr_service_limits}

\begin{table}[ht]
\centering
\caption{Whisper ба Gemini 3.5 Transcribe-ийн ялгаа}
\label{tab:asr_compare}
\begin{tabular}{|p{3.6cm}|p{5.0cm}|p{5.0cm}|}
\hline
 & \textbf{Whisper (Cafet small)} & \textbf{Gemini 3.5 Transcribe} \\
\hline
Ажиллах газар & Өөрийн компьютер & Google-ийн сервер \\
\hline
Интернэт & Хэрэггүй & Хэрэгтэй \\
\hline
Аудио & Тухайн компьютерт үлдэнэ & Google руу очно \\
\hline
Сайжруулах & Нэмэлт сургалт хийж болно; нэр томьёог эхний санамж болгон өгч болно \cite{openai_whisper_repo} & Загварыг сургах боломжгүй; нэр томьёоны жагсаалт өгч болно \\
\hline
Зардал & Компьютерын нөөц & Үнэгүй хэрэглээ, дараа нь төлбөртэй \\
\hline
Олон хүн зэрэг ашиглахад & Серверийн хүчин чадлаар хязгаарлагдана & Квот, төлбөрөөр хязгаарлагдана \\
\hline
\end{tabular}
\end{table}

\fullinput{Chapters/full/ch2_similar_systems_OFF}
